\section*{Capítulo 5: Estrellas de neutrones: las maravillas de la física}
\begin{description}[style=multiline,leftmargin=2em,labelwidth=1.5em]

% DOCX paragraph 2039
\item[1.]\hypertarget{note:05-01}{}Dos magnitudes fundamentales caracterizan un fenómeno ondulatorio, como el de las ondas electromagnética\index{radiación!electromagnética}s: la longitud de onda (es decir, la distancia entre dos «crestas» de la onda) y la frecuencia (es decir, el número de crestas que llegan en un intervalo de tiempo dado). La longitud de onda, \emph{λ}, y la frecuencia, \emph{f}, de la radiación electromagnética son inversamente proporcionales. Más precisamente, se relacionan como \emph{λ = c / f}, donde \emph{c} es la velocidad de la luz.

% DOCX paragraph 2040
\item[2.]\hypertarget{note:05-02}{}Este tipo de radiación\index{rayos X!radiación} también se llama \emph{rayos X blandos}, para distinguirla de la radiación\index{rayos X!radiación}\index{rayos X!radiación} de mayor energía, que se define como \emph{rayos X duros}. La dureza de un fotón en la banda de rayos X se mide, por tanto, mediante su energía (o, de forma equivalente, su frecuencia). Los rayos X blandos suelen tener una energía entre 0,1 y 0,3 keV (kiloelectronvoltios), mientras que los duros tienen entre 10 y 100 keV.

% DOCX paragraph 2041
\item[3.]\hypertarget{note:05-03}{}Conviene recordar que las primeras clasificaciones agrupaban los tipos de estrellas en \emph{«clases de colores»}, es decir, según el color dominante en el que aparecía su espectro\index{espectro!electromagnético} de emisión. Así, las estrellas se ordenaban desde el blanco azulado de las «más calientes» hasta el rojo anaranjado de las «más frías». En orden decreciente de temperatura o, de forma equivalente, de azul a rojo, las clases eran: O, B, A, F, G, K, M. Para ayudar a recordar la secuencia hay una frase sencilla introducida por la astrónoma estadounidense que trabajó en esta clasificación, Annie Jump Cannon (1863-1941): \emph{«Oh, be a fine girl/guy, kiss me!»}.

% DOCX paragraph 2042
\item[4.]\hypertarget{note:05-04}{}Para comprender bien por qué era un fenómeno tan extraño, convendría profundizar un poco en la cuestión. Es natural esperar que la escala temporal de variabilidad\index{variabilidad} de una fuente celeste sea proporcional a su tamaño. En particular, la luminosidad de los objetos grandes solo puede variar lentamente, mientras que la producida por objetos más pequeños (¡aunque a escala astronómica!) puede cambiar más deprisa. Esta relación entre tamaño y escala temporal es consecuencia de un requisito fundamental: la fuente está conectada causalmente. En otras palabras, sus cambios no pueden ser arbitrariamente rápidos, sino que deben producirse en una escala de tiempo (\emph{∆T}) mayor que el tiempo que tardaría un fotón, moviéndose a la velocidad de la luz (\emph{c}), en atravesar una fuente de un tamaño dado. Por eso, esta escala temporal se llama «tiempo de cruce de la luz». Matemáticamente, la restricción se traduce en una fórmula precisa: si \emph{R} es el tamaño característico de la fuente, para preservar la causalidad debemos imponer \emph{R \textless{} c ΔT}. Como \emph{∆T} se mide mediante las observaciones, obtenemos fácilmente un límite superior para el tamaño de la fuente, \emph{R}. Para darte una referencia, el tiempo de cruce de la luz en el Sol es de unos cinco segundos, mientras que su luminosidad está sujeta a variaciones que se producen en una escala temporal mucho más larga, de unos 11 años. Por tanto, ¡el Sol debe ser menor que 11 años luz, lo cual es obviamente correcto!

% DOCX paragraph 2043
\item[5.]\hypertarget{note:05-05}{}I. Shklovsky, \emph{Soviet Astronomy}, 11, 749 (1968)

% DOCX paragraph 2044
\item[6.]\hypertarget{note:05-06}{}W. Baade, F. Zwicky, \emph{Proceedings of the National Academy of Sciences of the United States of America}, 20, 254 (1934)

% DOCX paragraph 2045
\item[7.]\hypertarget{note:05-07}{}A. Hewish\index{Hewish, Antony}, S. J. Bell, J. D. H. Pilkington, P. F. Scott, R. A. Collins, \emph{Nature}, 217, 709 (1968)

% DOCX paragraph 2046
\item[8.]\hypertarget{note:05-08}{}La \emph{masa atómica} es el número que indica cuántos neutrones y protones contiene un determinado núcleo atómico. Una magnitud similar, pero distinta, es el \emph{número atómico}, que solo cuenta el número de protones de un núcleo atómico. Por ejemplo, el hidrógeno tiene masa atómica 1 y número atómico 1 porque su núcleo contiene un único protón; el helio tiene masa atómica 4 y número atómico 2 porque su núcleo contiene dos protones y dos neutrones.

% DOCX paragraph 2047
\item[9.]\hypertarget{note:05-09}{}Los neutrones son partículas fermiónicas, es decir, siguen una distribución estadística de Fermi--Dirac y, por tanto, están sujetos al principio de exclusión de Pauli. Este principio establece que un estado cuántico puede estar ocupado, como máximo, por una sola partícula de este tipo. Este mismo principio es responsable de la presión de degeneración\index{presión!degeneración}. Usando una analogía que, por desgracia, ha quedado demasiado clara en estos días, es como si las partículas fermiónicas estuvieran obligadas a mantener una «distancia social» y, por tanto, no pudieran estar demasiado cerca unas de otras. Cuando se acercan demasiado, generan mucha presión que impide una mayor compresión.

% DOCX paragraph 2048
\item[10.]\hypertarget{note:05-10}{}En mecánica cuántica\index{mecánica cuántica}, la función de onda representa el estado de un sistema físico, por ejemplo, de una partícula elemental. Es una función del tiempo y de la posición, de modo que se puede usar para medir la probabilidad de que una partícula se encuentre en una determinada región del espacio en un momento dado.

% DOCX paragraph 2049
\item[11.]\hypertarget{note:05-11}{}G. Gamow, M. Schoenberg, \emph{Physical Review}, 59, 539 (1941)

% DOCX paragraph 2050
\item[12.]\hypertarget{note:05-12}{}El modelo que propuse se llama «\emph{blitzar}» y plantea el nacimiento de un agujero negro a partir de una estrella de neutrones que colapsa y pierde su campo magnético. Las simulaciones numéricas de este proceso muestran que produciría una señal muy próxima a la observada, tanto en duración como en cantidad de energía emitida. Obviamente, esto no significa que el modelo sea correcto ni que explique realmente las observaciones: aún no tenemos el veredicto… (véase H. Falcke y L. Rezzolla, en \emph{Astronomy and Astrophysics}, 562, 134 (2014)).

% DOCX paragraph 2051
\item[13.]\hypertarget{note:05-13}{}F. Pacini, en \emph{Nature}, 219, 145 (1968)

% DOCX paragraph 2052
\item[14.]\hypertarget{note:05-14}{}He decidido presentar aquí un ejemplo distinto de PSR J0437-4715 para aclarar que estos no son púlsares «especiales».

\end{description}
